% Auto-generated by thesis/make_instruments.py --- do not edit by hand.
% Source: Novel Data/Headteacher Interview.docx

The schedule below is the instrument as administered. Each open-ended
question was scored 1--5 against the anchor descriptions shown. Each
question was asked by one researcher and scored by both; \cref{tab:irr_interview}
reports agreement between them, and the analysis uses the final recorded score
for each question. A checklist of key words and phrases was also completed for
each question; it enters no score reported in the chapter and is omitted here.
Item-to-sub-score mapping is given in \cref{app:2A:guide}.

\subsubsection*{Open-ended questions}

\paragraph{Q1. How long have you been in post?}
Recorded verbatim; not scored on the 1--5 scale.

\paragraph{Q2. How would you describe your school's culture?}
Rated 1 (weak answer) to 5 (strong answer).

\begin{description}
\item[1] A vague or generic description is provided, lacking specific details or examples.
\item[3] A detailed description of the school's culture, including some specific examples, but no explanation of the rationale for this culture.
\item[5] A detailed description of the school's culture with specific examples that also includes a clear explanation of the rationale for this culture.
\end{description}

\paragraph{Q3. How do you communicate this culture to your staff?}
Rated 1 (weak answer) to 5 (strong answer).

\begin{description}
\item[1] The strategies used to communicate the culture to the staff are not clearly explained.
\item[3] The strategies used to communicate the culture to the staff are clearly explained, but there is a lack of detail.
\item[5] There is a detailed and clear explanation of the strategies used to communicate the culture to the staff.
\end{description}

\paragraph{Q4. What do you think is the most important characteristic of an effective school? Why?}
Rated 1 (weak answer) to 5 (strong answer).

\begin{description}
\item[1] Various characteristics are mentioned but there is no clear explanation of why these characteristics lead to a school being effective.
\item[3] One or multiple characteristics are mentioned and there is a clear explanation of why these characteristics lead to a school being effective.
\item[5] One or multiple characteristics are mentioned and there is a clear explanation of why these characteristics lead to a school being effective. There is also a discussion of how the school culture links to these characteristics.
\end{description}

\paragraph{Q5. How does your reward and sanction system work? If a pupil misbehaves what should happen? If a pupil makes a positive contribution what should happen?}
Rated 1 (weak answer) to 5 (strong answer).

\begin{description}
\item[1] The explanation of the reward and sanction system is unclear, lacking specific details or a coherent structure.
\item[3] The response outlines the reward and sanction system to a decent extent, but there may be gaps in the explanation, and details might be limited.
\item[5] A clear and detailed explanation is provided on how the reward and sanction system operates; the response demonstrates a deep understanding of the procedures for addressing both misbehaviour and positive contributions.
\end{description}

\paragraph{Q6. What do you think are the key characteristics of an outstanding classroom teacher?}
Rated 1 (weak answer) to 5 (strong answer).

\begin{description}
\item[1] Various characteristics are mentioned but there is no clear explanation of why these characteristics make a teacher outstanding.
\item[3] Various characteristics are mentioned and there is a clear explanation of why these characteristics make a teacher is outstanding.
\item[5] Various characteristics are mentioned and there is a clear explanation of why these characteristics make a teacher is outstanding. There is also a discussion of how these characteristics link to the school's culture.
\end{description}

\paragraph{Q7. How important do you think extracurricular provision is? What kind of extracurricular provision have you put in place?}
Rated 1 (weak answer) to 5 (strong answer).

\begin{description}
\item[1] The importance of extracurricular provision is not clearly articulated, and the response lacks specific details about the types of activities available.
\item[3] The response acknowledges the importance of extracurricular provision, but the explanation may be somewhat basic or lacking in depth; details about the types of activities implemented could be more specific.
\item[5] The response shows a clear understanding of the importance of extracurricular provision and provides specific details about the types of activities in place.
\end{description}

\paragraph{Q8. What is your view on the importance of marking and feedback? Is there a school-wide policy on frequency and type?}
Rated 1 (weak answer) to 5 (strong answer).

\begin{description}
\item[1] The explanation of the importance of marking and feedback is unclear. If there is a school-wide policy, then it is also not explained clearly.
\item[3] The explanation of the importance of marking and feedback is clear. If there is a school-wide policy, then it is also explained clearly but without specific details.
\item[5] The answer lays out a clear and comprehensive view on the importance of marking and feedback. It provides specific details about a school-wide policy, including the frequency and type of marking and feedback.
\end{description}

\paragraph{Q9. What is the structure of your staff team?}
Rated 1 (weak answer) to 5 (strong answer).

\begin{description}
\item[1] The explanation is unclear and difficult to follow.
\item[3] There is a clear explanation of the staff structure from Senior Team through to Middle Leaders.
\item[5] The answer clearly lays out the staff structure from Senior Team down to classroom teachers and support staff with some explanation of the rationale for the structure.
\end{description}

\paragraph{Q10. How do you hold staff to account?}
Rated 1 (weak answer) to 5 (strong answer).

\begin{description}
\item[1] A brief answer which mentions the use of performance management structures but does not go into detail.
\item[3] A clear explanation of the performance management process with discussion of informal and formal steps.
\item[5] A clear explanation of the performance management process with discussion of informal and formal steps, as well as an explanation of an underlying philosophy.
\end{description}

\paragraph{Q11. How do you manage the wellbeing of your pupils?}
Rated 1 (weak answer) to 5 (strong answer).

\begin{description}
\item[1] An imprecise answer which emphasises the importance of managing pupil wellbeing but does not lay out any clear strategies for doing so.
\item[3] An answer which lays out strategies for managing pupil wellbeing, but some strategies are not very precise.
\item[5] An answer which lays out clear strategies for managing pupil wellbeing alongside an explanation of why this is important.
\end{description}

\paragraph{Q12. How do you assess pupil progress?}
Rated 1 (weak answer) to 5 (strong answer).

\begin{description}
\item[1] An answer which mentions no specific strategies for monitoring the progress that pupils are making.
\item[3] An answer which outlines strategies for assessing pupil progress in a summative manner only.
\item[5] An answer which outlines a range of strategies for assessing pupil progress formatively and summatively.
\end{description}

\subsubsection*{Yes-or-no questions}

Does your school make use of the following systems:

\begin{itemize}
\item Centralised detentions
\item An internal exclusion room
\item Reward assemblies/prize-giving ceremonies
\item Behaviour points/demerits
\item Reward points/merits
\item Line-ups
\item Silent corridors
\item After-school revision
\item Weekend revision
\item Ban on mobile phones in school
\end{itemize}

Would you like to say anything about your school's use of any of these systems?

\subsubsection*{Statement ranking}

For each of the following statements circle the appropriate number:

Each statement is rated 1 (strongly disagree) to 5 (strongly agree).

\begin{itemize}
\item The sanction systems are applied consistently by all staff
\item The reward systems are applied consistently by all staff
\item Teacher turnover is low
\item Staff morale is good
\item Pupil behaviour in class is generally good
\item Pupil behaviour outside of class is generally good
\item Pupils feel safe at school
\item Pupils feel like their teachers care
\item Teachers feel respected by all students
\item In general, parents are engaged with their child's education
\end{itemize}

Would you like to say anything about any of your responses to these questions?

Would you be willing to allow a few members of our team to come and observe the normal functioning of your school for a day?

\subsubsection*{To be completed after the interview}

\paragraph{Rate the confidence with which the headteacher responded to the questions}
Rated 1 (not confident at all) to 5 (confident).

\begin{description}
\item[1] There was a lot of hesitation throughout the interview and a lack of clarity in almost all the responses.
\item[3] The answers were generally clear and delivered without hesitation.
\item[5] Every answer was clear and delivered without hesitation. There were no wasted words.
\end{description}

\paragraph{Willingness to reveal information}
Rated 1 (very reluctant) to 5 (totally willing).

\begin{description}
\item[1] There was a reluctance to go into specific details throughout the interview.
\item[3] The headteacher went into specific detail with most answers, but not all.
\item[5] The headteacher was happy to go into very specific detail in every answer.
\end{description}

\paragraph{Interviewee patience}
Rated 1 (little patience) to 5 (lots of patience).

\begin{description}
\item[1] The headteacher rushed through all their answers.
\item[3] The headteacher gave detailed responses to most questions but did not add any detail to the yes/no or statement questions.
\item[5] The headteacher was happy to give detailed responses to all questions and wanted to add detail to their answers to the yes/no and statement questions.
\end{description}
